\documentclass[11pt]{article}
\usepackage[margin=0.9in]{geometry}
\usepackage{amsmath,amssymb,amsthm}
\usepackage{graphicx}
\usepackage[colorlinks=true,linkcolor=blue,urlcolor=blue,citecolor=blue]{hyperref}
\usepackage{microtype}

\newtheorem{theorem}{Theorem}
\newcommand{\vol}{\operatorname{vol}}
\newcommand{\R}{\mathbb R}

\title{A geometric block-polar proof of the mixed-ball volume factorization}
\author{Candidate full solution to a closing question of arXiv:1505.05867}
\date{11 August 2026}

\begin{document}
\maketitle

\noindent\textbf{Status.}
\texttt{candidate full solution (likely valid; human review requested)}.

\section*{Source problem}
Kempka and Vyb\'iral \cite{KV} prove an exact gamma-function formula for the
volume of the unit ball $B_{p,q}^{m,n}$ of $\ell_q^n(\ell_p^m)$.  In their
closing remarks (source PDF page~13), they observe that
\begin{equation}\label{source-factorization}
 \vol(B_{p,q}^{m,n})
 =\vol(B_q^{mn})
 \left(\frac{\vol(B_p^m)}{\vol(B_q^m)}\right)^n
\end{equation}
and ask whether this has a geometric or combinatorial interpretation giving a
non-analytic proof.  We give such an interpretation and prove a more general
star-body identity without evaluating any beta or gamma integral.

\section*{Transcription of the question}
Does formula \eqref{source-factorization} admit a geometric or combinatorial
interpretation which yields a non-analytic proof of the mixed-ball volume
formula?

\clearpage
\section*{Source evidence}
\noindent\textbf{The two closing questions and formula (19), source PDF
page 13.}
\begin{center}
\includegraphics[width=0.89\textwidth]{figures/source_page-13.png}
\end{center}

\clearpage
\section*{Geometric idea}
For a star body $K\subset\R^m$, a gauge shell at radius $r$ has volume
$m\vol(K)r^{m-1}\,dr$.  An outer $\ell_q$ constraint on $n$ blocks sees only
their gauge radii, so the admissible radial region is independent of the
angular shape of $K$.  Changing the inner body from $L$ to $K$ therefore
multiplies each of the $n$ shell measures by $\vol(K)/\vol(L)$.

Let $K$ be a bounded Borel star body containing the origin in its interior,
with gauge
\[
 \rho_K(x)=\inf\{t>0:x\in tK\}.
\]
For $0<q\le\infty$, define its $n$-fold $\ell_q$-sum by
\[
 \mathcal B_q^{(n)}(K)=
 \left\{(x_1,\ldots,x_n)\in(\R^m)^n:
 \|(\rho_K(x_1),\ldots,\rho_K(x_n))\|_{\ell_q^n}\le1\right\}.
\]

\begin{theorem}\label{ratio-theorem}
For any two such star bodies $K,L\subset\R^m$, any $n\ge1$, and
$0<q\le\infty$,
\[
 \frac{\vol_{mn}(\mathcal B_q^{(n)}(K))}
      {\vol_{mn}(\mathcal B_q^{(n)}(L))}
 =\left(\frac{\vol_m(K)}{\vol_m(L)}\right)^n.
\]
Convexity is unnecessary.
\end{theorem}

\begin{proof}
Homothety of gauge sublevels gives, for $r\ge0$,
\[
 \vol_m\{x:\rho_K(x)\le r\}=\vol_m(rK)=r^m\vol_m(K).
\]
Thus the pushforward of Lebesgue measure under $\rho_K$ is
\[
 (\rho_K)_\#(dx)=m\vol_m(K)r^{m-1}\,dr.
\]
Applying this identity independently in all $n$ blocks shows that, for every
nonnegative Borel function $F$,
\begin{align*}
 &\int_{(\R^m)^n}F(\rho_K(x_1),\ldots,\rho_K(x_n))\,dx_1\cdots dx_n\\
 &\qquad=(m\vol_m(K))^n
 \int_{[0,\infty)^n}F(r_1,\ldots,r_n)
 \prod_{j=1}^n r_j^{m-1}\,dr_1\cdots dr_n.
\end{align*}
Take $F$ to be the indicator of
\[
 R_{q,n}=\{r\in[0,\infty)^n:\|r\|_{\ell_q^n}\le1\}.
\]
Then
\[
 \vol_{mn}(\mathcal B_q^{(n)}(K))
 =(m\vol_m(K))^n I_{m,n,q},
 \qquad
 I_{m,n,q}=\int_{R_{q,n}}\prod_{j=1}^n r_j^{m-1}\,dr.
\]
The same constant occurs for $L$, so division proves the identity.  When
$q=\infty$, the identical argument uses $R_{\infty,n}=[0,1]^n$.
\end{proof}

\section*{A direct geometric transport}
The cancellation above can be realized by an actual map.  For a centrally
symmetric star body $K$, write
$S_K=\{\theta:\rho_K(\theta)=1\}$ and let $\mu_K$ be normalized cone measure.
The shell formula is the disintegration
\[
 dx=m\vol(K)r^{m-1}\,dr\,d\mu_K(\theta),\qquad x=r\theta.
\]
For $m\ge2$, the nonatomic standard probability spaces
$(S_L,\mu_L)$ and $(S_K,\mu_K)$ admit a measure-preserving bijection
$\phi:S_L\to S_K$ modulo null sets.  (For $m=1$, match the two endpoints.)
Its homogeneous extension
\[
 T(r\theta)=r\phi(\theta)
\]
preserves gauge radius, $\rho_K(Tx)=\rho_L(x)$, and scales volume by
$\vol(K)/\vol(L)$.  Hence $T^n$ maps
$\mathcal B_q^{(n)}(L)$ onto $\mathcal B_q^{(n)}(K)$ modulo null sets while
scaling volume by $(\vol(K)/\vol(L))^n$.  This gives a literal geometric
transport behind the factorization.

\section*{Resolution of the source question}
Take $K=B_p^m$ and $L=B_q^m$.  Then
\[
 \mathcal B_q^{(n)}(B_p^m)=B_{p,q}^{m,n},
 \qquad
 \mathcal B_q^{(n)}(B_q^m)=B_q^{mn}.
\]
Theorem~\ref{ratio-theorem} is precisely
\eqref{source-factorization}.  Geometrically, both bodies have the same
weighted block-radius base $R_{q,n}$; replacing every inner $q$-ball by an
inner $p$-ball changes only the $n$ angular shell factors.  This proves the
factorization without calculating the common radial integral.

The argument includes all $0<p,q\le\infty$: neither the triangle inequality
nor convexity is used, only positive homogeneity.  At $q=\infty$ it reduces to
the immediate product identity
$\vol((B_p^m)^n)=\vol(B_p^m)^n$.

\paragraph{Verification and novelty bounds.}
The proof uses only $\{\rho_K\le r\}=rK$, scaling of Lebesgue measure, and
Tonelli's theorem.  Checks at $n=1$, $m=1$, $p=q$, and $q=\infty$ give the
expected identities.  No formula for $I_{m,n,q}$ is used, so there is no
circular appeal to the source theorem.

A bounded search on 11 August 2026 covered the exact question and formula,
the source's citations, and combinations of ``geometric interpretation'',
``mixed-norm unit ball'', ``volume ratio'', ``block polar'', and ``cone
measure''.  No later source giving this exact shell proof was found.  The
source's separate Lorentz-volume question was answered later by
Dole\v{z}alov\'a and Vyb\'iral \cite{DV}; it is not claimed here.  A negative
search supports but cannot establish novelty.

\paragraph{Human-review recommendation.}
Verify the gauge-pushforward identity with the preferred harmless boundary
convention, and confirm that the source's requested ``non-analytical'' proof
includes a gamma-free geometric shell disintegration.  The volume identity
itself then follows immediately.

\small
\begin{thebibliography}{9}
\bibitem{KV}
H.~Kempka and J.~Vyb\'iral,
\emph{Volumes of unit balls of mixed sequence spaces},
arXiv:1505.05867; \href{https://arxiv.org/abs/1505.05867}{arXiv record}.

\bibitem{DV}
A.~Dole\v{z}alov\'a and J.~Vyb\'iral,
\emph{On the volume of unit balls of finite-dimensional Lorentz spaces},
arXiv:1906.04997; \href{https://arxiv.org/abs/1906.04997}{arXiv record}.
\end{thebibliography}

\end{document}
